% BEGIN THERMAL_R03_SYNC_13B
% Formalización reunida de resultados ya expuestos en el cuerpo.
\subsection{Definiciones estructurales y resultados cuantitativos}
\label{sec:ii-conclusiones-cuantitativas}

La definición por estructura conserva el objeto anterior a la lectura
escalar. Para las constantes de acción, ese objeto es una sección orientada
con una norma multiplicativa sobre el cociente local y un carácter regional.
Para Barbero--Immirzi, es
la evaluación de una incidencia orientada. Para la mezcla, es un transporte
entre coordenadas y amplitudes. Para Boltzmann, es una aplicación entre
rectas de magnitudes. Para la gravitación, es una sección recíproca de la
acción cuya magnitud procede de la longitud inscrita y de la lectura
radial especificadas en R1--R8. Estas definiciones
determinan operaciones diferentes dentro de la misma genealogía.

El cuadro siguiente fija, para cada magnitud, el objeto que la define,
la operación que determina su lectura y la función física desarrollada
en el cuerpo. Las pruebas de reconstrucción y transporte establecen
cómo se conserva esa información al reutilizar el resultado. En
particular, \ref{prop:registro-36-residuos} caracteriza la representación
del registro, \ref{mem:prop-schur} conserva la observación bajo reducción
y \ref{cor:ii-witt-composed-reader} conserva las coordenadas angulares
en la realización incidencial. Estas tres operaciones tienen dominios
propios y funciones demostrativas distintas dentro de la misma secuencia.

\begingroup
\setlength{\LTpre}{0pt}
\setlength{\LTpost}{2pt}
\begin{longtable}{@{}p{.13\linewidth}p{.415\linewidth}p{.395\linewidth}@{}}
\caption{Objeto estructural, lectura producida y significado físico de
los resultados del artículo.}\label{tab:ii-conclusiones-estructura}\\
\toprule
Magnitud & Definición estructural y resultado & Lectura física y desarrollo interno\\
\midrule
\endfirsthead
\toprule
Magnitud & Definición estructural y resultado & Lectura física y desarrollo interno\\
\midrule
\endhead
\bottomrule
\endfoot
Planck & Sección de la recta de acción:
$\hbar_{\rm pre}=H_5\,10^{-34}\mathcal U_S$,
$\hbar_{\rm ret}=\eta_{\rm ret}\,10^{-34}\mathcal U_S$.
El cociente local tiene dieciséis hojas y
$\operatorname{ord}_{10}(\varphi\alpha^{16})=-34$.
& La sección relaciona frecuencia y energía; la vuelta completa determina
$h_a=2\pi\hbar_a$. Desarrollo en \ref{sec:action}; coordenadas y comparación
en \ref{tab:ii-planck-comparacion}.\\[5pt]
Barbero--Immirzi & Evaluación orientada de los canales y de la incidencia:
$\gamma_{\HMT}=\pi AC^*/180+12\Delta S_{90}-\Delta S_{120}$.
La evaluación impresa es $0.2740076726944592747\ldots$.
& El coeficiente interviene en la resolución elemental de área y en
la realización Cartan--Holst. Las definiciones de
\ref{sec:barbero} y \ref{sec:area} conservan estos dos usos;
\eqref{eq:holst} especifica el segundo.\\[5pt]
Mezcla CKM & Imagen del vector angular mediante las tres aplicaciones
sectoriales; fase $\delta_{\deg}=9A_{\deg}$.
Los ángulos y amplitudes se evalúan después de construir el lector.
& El cuarteto de Jarlskog expresa la información de fase invariante
bajo cambios diagonales de base. Las reglas, su evaluación y la
realización unitaria se desarrollan en
\ref{sec:ii-propagacion-angular}.\\[5pt]
Boltzmann & Transducción positiva
$k_B:\mathcal L_\Theta\to\mathcal L_E$ con
$k_B\Theta_{{\rm clk},a}\log3=h_a/T_{108}$.
El estado y la medida de las continuaciones determinan el contenido
informacional sobre el que actúa.
& $S=k_Bs$ realiza dimensionalmente la entropía del estado reducido.
Las bases de energía y temperatura y sus transformaciones se conservan
en \ref{sec:ii-boltzmann}; la caracterización aditiva y el instrumento
se demuestran en \ref{sec:ii-aditiva-instrumento}; la comparación está en
\ref{tab:ii-kb-g-cartas}.\\[5pt]
Realización térmica marcada & Referencia fija $\eta$, con multiplicidades
$(d_0,d_1,d_2)=(1,380,649)$ en grados $23,26,29$ y traza
$b_*=1.380649\times10^{-23}$.
Las lecturas $\mathcal S_P=b_*s$ y $\mathcal E_P=\operatorname{Tr}(\rho H)$
determinan $d\mathcal S_P/d\mathcal E_P=b_*\beta$ y
$\Theta_P=(b_*\beta)^{-1}$.
& Con referencia fija, bases positivas transportadas y Hamiltoniano no
escalar, \ref{thm:ii-ter27-transducer} determina el transductor.
El corolario~\ref{cor:ii-ter27-variational} prueba el mínimo único.
La composición conserva $q^N/Z_N(q)$, el origen energético y los
transportes reloj--horizonte de \eqref{eq:ii-ter27-intertwiner}.\\[5pt]
Gravitación & Sección recíproca
$G_a=c_{\rm int}^{3}L^2/\hbar_a$,
con $L=(5759/23040)\alpha_{\HMT}^{16}L_*$.
El reloj $t_0=\pi_{\HMT}L/(54c_{\rm int})$ conserva la expresión
circular $G_a=\vartheta_{108}^{\circ}c_{\rm int}^{5}t_0^2/\hbar_a$,
con $\vartheta_{108}^{\circ}=(54/\pi_{\HMT})^2$.
& El producto $G_a\hbar_a$ conserva la normalización areal.
La construcción longitudinal y radial, sus reglas y la unicidad del
coeficiente se demuestran en \ref{g26:sec:rigidez-radial-es};
la realización circular se conserva en \ref{sec:gravity} y la confluencia térmica en
\ref{sec:confluencia}.\\
\end{longtable}
\endgroup

La comparación cuantitativa conserva el significado propio de cada
referencia. En Planck, la tabla~\ref{tab:ii-planck-comparacion} compara
coordenadas dimensionales con una constante definitoria del SI; los
residuos relativos de las orientaciones anterior y posterior son,
respectivamente, aproximadamente $-1.82\times10^{-15}$ y
$-6.93\times10^{-10}$. En Barbero--Immirzi, la
tabla~\ref{tab:ii-barbero-comparacion} compara dos construcciones:
la evaluación interna y la raíz de la ecuación de Ghosh--Mitra, cuyo
valor aproximado es $0.2740668582328300$. La diferencia relativa es
aproximadamente $-2.16\times10^{-4}$; el objeto comparado es un
coeficiente teórico.

En CKM, la tabla~\ref{tab:ii-ckm-comparacion} reúne los senos de los
tres ángulos de mezcla, la fase en radianes y el invariante de Jarlskog.
Los ángulos y la matriz de módulos se explicitan en el desarrollo
contiguo.
Las desviaciones marginales normalizadas allí definidas satisfacen
$|z|<0.029$ frente a las filas PDG utilizadas. Cada residuo utiliza
la incertidumbre marginal indicada. Su lectura es individual;
una evaluación conjunta requiere además la matriz de covarianza
de esos observables.
Para Boltzmann y gravitación, la tabla~\ref{tab:ii-kb-g-cartas}
expone el transductor, la composición longitudinal y radial, las coordenatizaciones dimensionales
y la referencia externa por separado. Esta organización permite
identificar exactamente qué cantidad produce cada construcción y
qué cantidad evalúa cada contraste.

El cuadro~\ref{tab:ii-conclusiones-contraste} reúne los resultados de
esas comparaciones junto a las definiciones estructurales anteriores.
\begin{table}[!htbp]
\centering\small
\setlength{\tabcolsep}{4pt}
\renewcommand{\arraystretch}{1.12}
\begin{tabularx}{\linewidth}{@{}>{\raggedright\arraybackslash}p{.16\linewidth}>{\raggedright\arraybackslash}X>{\raggedright\arraybackslash}X@{}}
\toprule
Magnitud & Resultado en la realización especificada & Referencia y contraste\\
\midrule
Planck, anterior al retorno
& \(\hbar_{\rm pre}=1.0545718176461544696\ldots\)
  \(\times10^{-34}\,\mathrm{J\,s}\).
& SI: \(\hbar_{\rm SI}=h_{\rm SI}/(2\pi)\), exacto.
  Su coordenada es \(1.0545718176461563912\ldots\)
  \(\times10^{-34}\,\mathrm{J\,s}\).
  Residuo relativo \(-1.822221\times10^{-15}\).\\[4pt]
Planck, posterior al retorno
& \(\hbar_{\rm ret}=1.0545718169150390611\ldots\)
  \(\times10^{-34}\,\mathrm{J\,s}\).
& Misma referencia SI.
  Residuo relativo \(-6.932836\times10^{-10}\).
  Para \(h_a=2\pi\hbar_a\), los residuos son idénticos.\\[4pt]
Barbero--Immirzi
& \(\gamma_{\HMT}=0.2740076726944592747\ldots\).
& Ghosh--Mitra: aproximadamente \(0.2740668582328300\).
  Diferencia relativa \(-2.159529202\times10^{-4}\)
  entre construcciones teóricas.\\[4pt]
Mezcla CKM
& \(\delta_{\rm CP}=65.676173123554^\circ\),
  \(J=3.1189723326632974\times10^{-5}\).
  Los tres ángulos de \eqref{eq:ii-ckm-decimales} son
  \(\theta_{12}=13.003313589509^\circ\),
  \(\theta_{23}=2.398056157332^\circ\) y
  \(\theta_{13}=0.213977382244^\circ\).
& PDG 2025: las cinco comparaciones de
  \(s_{12},s_{23},s_{13},\delta,J\) satisfacen \(|z_x|<0.029\)
  con las incertidumbres marginales de la
  tabla~\ref{tab:ii-ckm-comparacion}.\\[4pt]
Boltzmann
& \(B_*=1.380649\times10^{-23}\,u_E/u_\Theta\),
  con las bases y la regla termométrica de
  \eqref{eq:ii-b-evaluacion-comparada}.
& SI: \(1.380649\times10^{-23}\,\mathrm{J/K}\), exacto.
  Coincidencia del numeral bajo el transporte dimensional declarado.\\[4pt]
Longitud \(L\)
& \(L=1.6162549826251263301\ldots\)
  \(\times10^{-35}\,\mathrm m\), para \(L_*=1\,\mathrm m\).
& Composición longitudinal y radial:
  \(G_{\rm ret}=c_{\rm int}^3L^2/\hbar_{\rm ret}\),
  con las mismas bases de velocidad y acción.\\[4pt]
Gravitación
& \(G_{\rm ret}=6.6742996594140227\ldots\)
  \(\times10^{-11}\,\mathrm{m^3\,kg^{-1}\,s^{-2}}\).
& CODATA 2022: \(6.67430(15)\times10^{-11}\) en las mismas unidades.
  Diferencia de aproximadamente \(-0.00227057\) incertidumbres estándar.\\
\bottomrule
\end{tabularx}
\caption{Síntesis del contraste cuantitativo. Las evaluaciones, las bases
y las referencias son las de las tablas~\ref{tab:ii-planck-comparacion}--\ref{tab:ii-kb-g-cartas}.
Los residuos relativos de las constantes de acción, la comparación
teórica de Barbero--Immirzi y los residuos normalizados de CKM y gravitación
conservan sus respectivos significados.}
\label{tab:ii-conclusiones-contraste}
\par\medskip
\centering\small
\setlength{\tabcolsep}{4pt}
\renewcommand{\arraystretch}{1.18}
\begin{tabular}{@{}lrrr@{}}
\toprule
Observable & Evaluación HMT & PDG 2025 & \(z_x\)\\
\midrule
\(s_{12}\) & \(0.225007404757\) & \(0.22501\pm0.00068\) & \(-0.003817\)\\
\(s_{23}\) & \(0.041841757010\) & \(0.04183^{+0.00079}_{-0.00069}\) & \(0.014882\)\\
\(s_{13}\) & \(0.003734601164\) & \(0.003732^{+0.000090}_{-0.000085}\) & \(0.028902\)\\
\(\delta\) (rad) & \(1.146265461116\) & \(1.147\pm0.026\) & \(-0.028251\)\\
\(10^5J\) & \(3.118972333\) & \(3.12^{+0.13}_{-0.12}\) & \(-0.008564\)\\
\bottomrule
\end{tabular}
\caption{Valores CKM y comparación marginal reunidos en las conclusiones.
Los tres ángulos de la tabla~\ref{tab:ii-conclusiones-contraste} se contrastan
mediante \(s_{ij}=\sin\theta_{ij}\), junto con la fase en radianes y el
invariante \(J\). Se conservan las cifras, incertidumbres y residuos de la
tabla~\ref{tab:ii-ckm-comparacion}, correspondientes al ajuste de tres
generaciones PDG 2025. La última fila expresa conjuntamente valor e
incertidumbre en unidades de \(10^{-5}\). Los residuos son comparaciones
marginales de observables correlacionados.}
\label{tab:ii-conclusiones-ckm}
\end{table}


\Needspace{10\baselineskip}
\subsection{Conservación de la escala areal y recuperación de la información de frontera}

La relación entre información y gravitación se concentra en dos mapas
complementarios. El primero conserva la escala areal mediante
$G_a\hbar_a$ y organiza sus ocupaciones mediante
$\gamma_{\HMT}a_0$. El segundo conserva la resolución sectorial mediante
la observación conjunta del área y de la incidencia ponderada. La inversión
\eqref{eq:recover} recupera las ocupaciones desde el área y el
observable ponderado; las ecuaciones
\eqref{eq:ii-respuesta-entropica} y \eqref{eq:ii-respuesta-areal}
vinculan sus variaciones con las fluctuaciones del estado.

De aquí procede la lectura física central del artículo: una misma
configuración admite una observación geométrica y una observación
informacional relacionadas por una estructura recuperable. El cambio
de orientación puede conservar la entropía del estado reducido y
transformar el área, como establecen
\eqref{eq:ii-inversion-entropia} y \eqref{eq:ii-area-complementaria}.
La conservación de información se refiere al estado y a su transporte;
la entropía de entrelazamiento se refiere además a la bipartición y a
la reducción. La identidad de horizonte de
\ref{sec:confluencia} utiliza esas funciones en la realización
geométrica y térmica especificada.

Las identidades de memoria permiten precisar también el alcance de una
reducción de variables. La igualdad de observaciones de
\eqref{mem:schur-lector} conserva las contribuciones del sector eliminado
mediante el operador efectivo, la fuente y el lector transformados.
La entropía del estado reducido, en cambio, se evalúa con la bipartición
y la medida especificadas en \ref{sec:ii-ecuacion-informacion}.
Ambas construcciones explican funciones diferentes de la memoria:
la primera conserva una lectura del transporte; la segunda cuantifica
la información accesible al subsistema. Su articulación mantiene
explícito qué estructura se transporta y qué reducción define el
observable físico.

El avance explicativo se expresa, por tanto, mediante un orden preciso:
estructura generada, aplicación de lectura, valor obtenido y función
física. La igualdad de dos valores se estudia junto con la igualdad
o diferencia de sus estructuras; los mapas de recuperación determinan
cuál de esas diferencias conserva significado observable.
% END THERMAL_R03_SYNC_13B
